% ========== 
% AESTHETICS 
% ==========

\usetheme[progressbar=frametitle]{metropolis}
\usepackage{appendixnumberbeamer}
\usepackage{booktabs}
\usepackage[scale=2]{ccicons}
\usepackage{pgfplots}
\usepgfplotslibrary{dateplot}
\usepackage{xspace}
\newcommand{\themename}{\textbf{\textsc{metropolis}}\xspace}

% 1. Define the official UT Burnt Orange
\definecolor{UTBurntOrange}{RGB}{191, 87, 0}
\definecolor{UTCharcoal}{RGB}{51, 63, 72} % A great secondary color for text/titles

% 2. Apply the color to Beamer elements
\setbeamercolor{structure}{fg=UTBurntOrange}        % Changes bullets, numbers, etc.
\setbeamercolor{title}{fg=UTBurntOrange}            % Changes the main title
% \setbeamercolor{frametitle}{fg=UTBurntOrange}       % Changes the header of each slide
\setbeamercolor{palette primary}{bg=white, fg=UTBurntOrange} % For bars/footers
\setbeamertemplate{progress bar in head/foot}{}
\setbeamercolor{background canvas}{bg=white}

\setbeamercolor{block title}{bg=UTBurntOrange!10, fg=gray}
\setbeamercolor{block body}{bg=UTBurntOrange!10, fg=black} % 10% transparency for the body

\setcounter{tocdepth}{1}


% ========== 
% Additional macros 
% ========== 

\usepackage{algorithmic}
\usepackage{algorithm}
% \usepackage{algpseudocode}
\renewcommand{\figurename}{Fig.}
\usepackage[font=footnotesize, labelfont={bf, color=structure}]{caption}

\usepackage{setspace}

\usepackage{xcolor}

\usepackage{tikz}
\usetikzlibrary{arrows.meta}

\usepackage{multirow}
